{\Large\textbf{Ideal Gas, Vapor Pressure, and Heat Conduction [20 pts.]}}

In this experimental exam, you will study three physical phenomena: the
behavior of an ideal gas, the vapor pressure of water, and heat transfer by
conduction. The objective of the exam is to analyze the relationships between
pressure $P$, air volume $V$, temperature $T$, and time $t$, present your
results in relevant graphs, and estimate associated physical constants and
their associated uncertainties.

\begin{center}\includegraphics[width=0.95\linewidth]{figures/IPhO_2026_Q4_fig1.jpg}\end{center}

{\large\textbf{Overview of the Experimental Setup}}

\textbf{Before beginning the experiment, read the assembly description, list
of materials, general instructions, and procedure for each part of the test.}

The apparatus consists of two concentric cylinders made of acrylic (Fig. 1a).
The inner cylinder (hereafter \textbf{IC}) can contain air and liquids, and the
air inside \textbf{IC} is referred to as confined air (hereafter \textbf{CA}).
The amount of liquid in the \textbf{IC} can be controlled by connecting the
syringe to \textbf{hose D}.\\
The outer cylinder (hereafter \textbf{OC}) can only be filled with water and
acts as a thermal jacket. Its water level can be adjusted using inlet and
outlet hoses and valves, while its temperature is controlled by an electric
heater. Two sensors are installed inside the \textbf{IC}: one for pressure and
one for temperature. A separate temperature sensor is available in the
\textbf{OC}.

\begin{center}\includegraphics[width=0.8\linewidth]{figures/IPhO_2026_Q4_fig2.jpg}\end{center}

{\large\textbf{Experimental Equipment}}

\textbf{\textit{Main Apparatus (see Fig. 2)}}

1. \textbf{Inner cylinder} (\textbf{IC}): Graduated with two centimeter scales.
Its fixed lids keep the contents airtight.

2. \textbf{Outer cylinder} (\textbf{OC}): Acts as a thermal jacket.

3. \textbf{Heater:} Installed in the \textbf{OC} and powered by the electronic
module.

4. \textbf{Temperature sensors}: Two sensors, connected to the electronic
module by black wires, measure the temperature in each cylinder.

5. \textbf{Pressure sensor hose:} Connects \textbf{CA} to a pressure sensor
hosted in the electronic module, allowing measurement of the internal
pressure.

6. \textbf{OC water connections}: Two hoses provide water with an inlet to and
outlet from the \textbf{OC.}

\textbf{Valve E — IC atmospheric connection:} Opens or closes the connection
between the \textbf{IC} and the atmosphere.

\textbf{\textit{Electronic Module (see Fig. 3)}}

The electronic module must be connected to the main power before use. The
power switch, labeled \textbf{Power}, is located at the back of the module
(Fig. 4).

The digital display provides the following readings:

\textbullet\ \textbf{Internal pressure:} Pressure inside the \textbf{IC}
(\textbf{INT. PRESSURE}).

\textbullet\ \textbf{Atmospheric pressure:} Local atmospheric pressure
(\textbf{ATM. PRESSURE}).

\textbullet\ \textbf{Internal temperature:} Temperature inside the \textbf{IC}
(\textbf{INT. TEMPERATURE}).

\textbullet\ \textbf{External temperature:} Temperature inside the \textbf{OC}
(\textbf{EXT. TEMPERATURE}).

\textbullet\ \textbf{Stopwatch:} Timer reading (\textbf{TIMER}).

Four numbered buttons below the screen control the main module functions:

\textbullet\ \textbf{Button 1} — \textbf{Timer}: Starts or stops the
stopwatch.

\textbullet\ \textbf{Button 2} — \textbf{Reset}: Clears the timer reading.

\textbullet\ \textbf{Button 3} — \textbf{Graph display}: toggles between Graph
display and numerical display. \textbf{This function is not used in the
experiments.}

\textbullet\ \textbf{Button 4} — \textbf{Heater}: Turns the heater on or off.
When the heater is active, the animated indicator in the lower-right corner of
the screen alternates colors between red and green.

\textbf{Note: Never turn on the heater if it is not fully immersed in water.
When the OC temperature reaches 65\,°C, the heater is deactivated. It may be
reactivated only after the animated indicator is no longer red.}

The right side of the module (Fig. 5) contains the connections for the
\textbf{IC} temperature sensor (Internal Temp.), the \textbf{OC} temperature
sensor (\textbf{External Temperature}), the heater (\textbf{Water Heater}), and
the hose that connects \textbf{CA} to the pressure sensor (\textbf{Internal
Pressure})

\textbf{\textit{Hydraulic Module (see Fig. 6)}}

The hydraulic module contains four valves, labeled \textbf{A}, \textbf{B},
\textbf{C}, and \textbf{D}. Each valve controls a specific fluid pathway in the
apparatus:

\textbullet\ \textbf{Valve A — Water inlet to OC}: Controls the flow of water
into the \textbf{OC} through \textbf{hose A} (\textbf{IN H2O}).

\textbullet\ \textbf{Valve B — Recirculator valve:} Controls the flow of water
through the recirculation cycle, described in the pump operation instructions
below.

\textbullet\ \textbf{Valve C — Water outlet from OC:} Controls the flow of
water out of the \textbf{OC} through \textbf{hose C} (\textbf{OUT H2O}).

\textbullet\ \textbf{Valve D — Syringe connection:} Controls the inlet and
outlet of fluids in the \textbf{IC} using the syringe connected to \textbf{hose
D} (Fig. 2). If the valve is leaky you can detach the hose from the setup and
connect the syringe directly to the hose.

\textbf{\textit{Pump and Valve Operation}}

The pump switch (\textbf{PUMP}), located on the front of the electronic module
(Fig. 8), activates the pump inside the electronic module. Before turning on
the pump, set the valves according to the operation you want to perform. Use
the clips to secure the hoses to the edge of the containers before operating
the pump.

\textbullet\ \textbf{Fill the OC with water:} Open \textbf{valve A} and close
\textbf{valves B} and \textbf{C}. Place \textbf{hose A} in a water container.
When the desired level is reached, close \textbf{valve A}. Make sure the water
level remains at least 1\,cm below the cover of the cilinder. \textit{If water
does not enter \textbf{OC} efficiently, immediately open and close
\textbf{valve B} to purge the pump.}

\textbullet\ \textbf{Drain water from the OC:} Close \textbf{valve A} and open
\textbf{valves B} and \textbf{C}. Place \textbf{hose C} in a container to
collect the water.

\textbullet\ \textbf{Homogenize the OC temperature:} Open \textbf{valve B} and
close \textbf{valves A} and \textbf{C}. In this configuration, water circulates
through the \textbf{OC} in a closed loop, helping to make the temperature
uniform.

For the following operations, connect the syringe to \textbf{hose D}. The pump
is not used in these cases.

\textbullet\ \textbf{Add liquid to the IC:} Open \textbf{valves D} and
\textbf{E,} then use the syringe to introduce the liquid.

\textbullet\ \textbf{Remove liquid from the IC:} Open \textbf{valves D} and
\textbf{E}, then use the syringe to draw out the liquid.

\begin{center}\includegraphics[width=0.9\linewidth]{figures/IPhO_2026_Q4_fig3.png}\end{center}

\textbf{Additional Items: the following materials are supplied for the
experiment.}

\textbullet\ \textbf{Two 100\,mL syringes} with a connecting hose (Fig. 7).
One syringe is for water, the other for \textbf{PG}.

\textbullet\ \textbf{Propylene glycol (PG):} A container with 100\,mL of
\textbf{PG}, whose vapor pressure is negligible over the temperature range used
in this experiment (Fig. 9). \textbf{Warning: PG is miscible with water.} It
can also hydrate if open to the air. It is recommended that you keep the
container closed while not using it.

\textbullet\ \textbf{Cold water:} A thermal container with 500\,mL of cold
water (Fig. 10).

\textbullet\ \textbf{Room-temperature water:} A bottle of water at room
temperature (Fig. 11).

\textbullet\ \textbf{Pouring container:} A container used for pouring water
(Fig. 12).

\textbullet\ \textbf{Excess-water jug:} A jug for collecting excess water
(Fig. 13).

\textbullet\ \textbf{Clips:} For securing hoses and supporting components
during the experiment (Fig. 14).

\textbullet\ \textbf{Paper towels} for cleanup.

\textbf{Important Operating Notes}

Keep the following notes in mind when setting up and operating the apparatus:

\textbullet\ \textbf{Valve positions:} Each valve has only two positions: open
(\textbf{I}) and closed (\textbf{II}) (Fig. 15).

\textbullet\ \textbf{Hose connection:} To remove a hose, push the hoop and pull
the hose out. To insert a hose, push it fully into place (Fig. 16).
\textbf{Note that the apparatus is fully assembled at the start of the exam and
the hoses normally do not need to be reconnected in the usual operation of the
apparatus.}

\textbullet\ \textbf{Cylinder dimensions:} Figure 17 shows the cross-sectional
dimensions of the cylinders.

\textbullet\ \textbf{Notation:} Figure 18 defines the notation used for the
\textbf{IC} contents: \textbf{\textit{h}} is the height of the liquid, and
\textbf{\textit{H}} is the height of the air above the liquid.

\textbullet\ \textbf{Volume corrections:} The internal volume of the connecting
hoses is approximately compensated by the volume of the temperature sensor
inside the \textbf{IC}. Therefore, these effects should be neglected.

\textbf{\textit{Constants and Reference Values}}

The following constants and reference values are used throughout the
experiment:

\textbullet\ \textbf{Avogadro constant:} $N_A = 6.022 \times 10^{23}\,\mathrm{mol}^{-1}$

\textbullet\ \textbf{Specific heat capacity of water:} $c_{\mathit{0}} = 4.184 \times 10^{3}\,\mathrm{J/(kg{\cdot}K)}$

\textbullet\ \textbf{Density of water:} $\rho_{\mathit{0}} = 1000\,\mathrm{kg/m^3}$

\textbullet\ \textbf{Molar mass of water:} M\textsubscript{\textit{0}} = 18.02\,g/mol

\textbullet\ \textbf{Molar mass of air:} M\textsubscript{\textit{a}} = 28.96\,g/mol

\textbullet\ \textbf{Molar volume of an ideal gas at normal conditions:}
$V_m = 22.4\,\mathrm{L/mol}$, for $T_{\mathit{0}} = 273.15\,\mathrm{K}$ and
$P_{\mathit{0}} = 101.3\,\mathrm{kPa}$

{\large\textbf{Part A [4.0 pt]. Isochoric (isovolumetric) process of an ideal gas}}

\textbf{CA} obeys the ideal gas equation of state,
\begin{equation}
PV = nRT \tag{1}
\end{equation}
where $R$ is the universal gas constant.

\textbf{Procedure}

Introduce \textbf{PG} into \textbf{IC} to $h = 4.5\,\mathrm{cm}$ and close
valves \textbf{D} and \textbf{E.} This ensures that the volume of \textbf{CA}
is fixed.

The density of ambient air in Bucaramanga varies with local temperature and
pressure. For the following question, use the time-averaged value
$\uprho = 1.12\,\mathrm{kg/m^3}$.

\noindent\fbox{\parbox{\dimexpr\linewidth-2\fboxsep-2\fboxrule\relax}{\textbf{A1.}
[0.4 pt] Determine the mass $m$, number of moles $n$, and the total number of
air molecules $N$ of the \textbf{CA}.}}

1. Fill \textbf{OC} with room temperature water close to the top.

2. Turn on the heater. Don't forget to homogenize the temperature by using the
pump.

\noindent\fbox{\parbox{\dimexpr\linewidth-2\fboxsep-2\fboxrule\relax}{\textbf{A2.}
[ 1 pt] Record in a table the pressure $P$ of the CA as a function of its
temperature T.}}

\noindent\fbox{\parbox{\dimexpr\linewidth-2\fboxsep-2\fboxrule\relax}{\textbf{A3.}
[0.9pt] From the data obtained in question A2, plot the behavior of pressure as
a function of temperature.}}

\textbf{\textit{Under IPhO competitive conditions, the sensors may be
intentionally decalibrated to alter the expected value of R with respect to the
reference value.}}

\noindent\fbox{\parbox{\dimexpr\linewidth-2\fboxsep-2\fboxrule\relax}{\textbf{A4.}
[1.0pt] Use the graph from A.3 to determine the experimental value of $R$.}}

The thermal pressure coefficient at constant volume is defined by the
relationship:
\begin{equation}
\beta_0 = \tfrac{1}{P_0}\left(\tfrac{\Delta P}{\Delta T}\right), \tag{2}
\end{equation}
where $P_0$ is the pressure of the system at the reference temperature $T_0$
as indicated in the reference constants and values.

\noindent\fbox{\parbox{\dimexpr\linewidth-2\fboxsep-2\fboxrule\relax}{\textbf{A5.}
[0.7pt] Determine the value of the coefficient $\upbeta_0$ for air.}}

\textbf{\textit{Remove the PG from the IC and clean it before proceeding to the
other parts. If you proceed to part B, do not remove the hot water from OC.}}

{\large\textbf{Part B [8.0 pt]. Vapor Pressure}}

In a closed container containing a liquid and its vapor at equilibrium, the
vapor pressure depends on the absolute temperature T according to the
Clausius–Clapeyron equation:
\begin{equation}
P_v = P_{v0} e^{-\frac{Q_v}{R}\left(\frac{1}{T} - \frac{1}{T_0}\right)} \tag{3}
\end{equation}
Here, $T_{\mathit{0}}$ is the reference temperature, taken as 0\,°C.
$P_{v\mathit{0}}$ is the vapor pressure at $T_{\mathit{0}}$. $R$ is the
universal gas constant, and Q\textsubscript{\textit{v}} is the molar latent
heat of vaporization of water.

\begin{center}\includegraphics[width=0.25\linewidth]{figures/IPhO_2026_Q4_fig4.png}\end{center}

\textbf{Procedure}

\textbf{\textit{Remember that PG and water are miscible. Use different syringes
for water and PG}}

\textbullet\ Fill the \textbf{OC} with water until the water level is close to
the top.

\textbullet\ Fill a syringe with water, connect it to \textbf{IC} and open
valve \textbf{E}, remove the syringe plunger, and raise the syringe until it
reaches the water level at $h = 5.0\,\mathrm{cm}$. Close valve \textbf{E} and
then secure the syringe to the jug as shown in Figure 19. This arrangement
keeps the pressure inside the \textbf{IC} approximately equal to the
atmospheric pressure. Make sure no air bubbles remain in the hose.

\textbullet\ Heat the water in the \textbf{OC} to 65\,°C. Remember to use the
recirculator system to homogenize the water temperature.

\textbullet\ Wait until the temperature starts to decrease to start recording
$H$ as a function of $T$.

\textbullet\ To lower the temperature further, you may replace the water in the
\textbf{OC} with the supplied cold water at the appropriate time.

\noindent\fbox{\parbox{\dimexpr\linewidth-2\fboxsep-2\fboxrule\relax}{\textbf{B1.}
[1.0 pt] Record the height H as a function of temperature T in a table}}

\noindent\fbox{\parbox{\dimexpr\linewidth-2\fboxsep-2\fboxrule\relax}{\textbf{B2.}
[1.0 pt] Use the data obtained in B1 to plot $H$ as a function of $T$.}}

The vapor pressure is relatively small around room temperature and the
$H$($T$) curve may be approximated as linear.

\noindent\fbox{\parbox{\dimexpr\linewidth-2\fboxsep-2\fboxrule\relax}{\textbf{B3.}
[0.5 pt] Use the graph from question B2 to extrapolate and find the value of
H\textsubscript{0} ($H$ at 0\,°C).}}

\noindent\fbox{\parbox{\dimexpr\linewidth-2\fboxsep-2\fboxrule\relax}{\textbf{B4.}
[1.0 pt] Assuming that \textbf{IC} contains both dry air plus water vapor,
deduce an algebraic expression for the vapor pressure \textit{P}\textsubscript{v}
in terms of $P_{atm}$ \textit{(atmospheric pressure),} $H_{\mathit{0}}$,
$H$, $T_{\mathit{0}}$ and $T$. For the purposes of this problem, you can assume
that the vapor pressure at 0\,°C is zero.}}

For the following questions, use the reference value
$R = 8.31\,\mathrm{J/(mol{\cdot}K)}$.

\noindent\fbox{\parbox{\dimexpr\linewidth-2\fboxsep-2\fboxrule\relax}{\textbf{B5.}
[4.0 pt] Using equation (3) and the expression from question B4, construct an
appropriate graph to determine the experimental value of Q\textsubscript{v.}}}

\noindent\fbox{\parbox{\dimexpr\linewidth-2\fboxsep-2\fboxrule\relax}{\textbf{B6.}
[0.5 pt] From the value of Q\textsubscript{\textit{v}}, determine
$L_v$,\textit{(the latent heat of vaporization per unit mass)} indicating the
formula that you used.}}

{\large\textbf{Part C [8.0 pt]. Heat Conduction}}

In this part of the exam, you will study the heat exchange between the water in
the \textbf{OC} and the water in the \textbf{IC} through the wall that
separates them. By analyzing the time evolution of the water temperature in the
\textbf{IC}, $T_{IC}$, and the water temperature in the \textbf{OC}, $T_{OC}$,
the objective of this section is to estimate the effective thermal resistance
of the wall, the equilibrium temperature of the system, and the thermal
conductivity of acrylic.

Heat conduction through a solid wall follows the relation:
\begin{equation}
\tfrac{\Delta Q}{\Delta t} = \tfrac{1}{R_{Th}}(T_{OC} - T_{IC}) \tag{4}
\end{equation}
where $\Delta Q$ is the heat received by the water in the \textbf{IC} through
the wall during the time interval $\varDelta t$. The effective thermal
resistance $R_{Th}$ depends on the material and geometry of the wall separating
\textbf{IC} and \textbf{OC}.

\textbf{Procedure}

1. Set the water level in the \textbf{OC} to $h = 15\,\mathrm{cm}$.

2. Heat the water in the \textbf{OC} to 65\,°C. Use the pump to homogenize the
water temperature.

3. Set the water level in the \textbf{IC} to $h = 10\,\mathrm{cm}$ and start
the stopwatch.

\noindent\fbox{\parbox{\dimexpr\linewidth-2\fboxsep-2\fboxrule\relax}{\textbf{C1.}
[1.0 pt] Record the internal temperature $T_{IC}$ and external temperature
$T_{OC}$ as a function of time t}}

\noindent\fbox{\parbox{\dimexpr\linewidth-2\fboxsep-2\fboxrule\relax}{\textbf{C2.}
[1.0 pt] On a single graph, plot $T_{IC}$ and $T_{OC}$ as functions of time
$t$.}}

For questions C3 and C4, ignore the heat capacity of the apparatus.
T\textsubscript{eq} refers to the equilibrium temperature that would be reached
by both IC and OC if there were no heat transfer to the environment.

\noindent\fbox{\parbox{\dimexpr\linewidth-2\fboxsep-2\fboxrule\relax}{\textbf{C3.}
[0.7 pt] Plot T\textsubscript{OC} - T\textsubscript{IC} as functions of both
T\textsubscript{OC} and T\textsubscript{IC} on the same graph with
T\textsubscript{OC} and T\textsubscript{IC} together on the same temperature
axis and T\textsubscript{OC} - T\textsubscript{IC} on the other axis.}}

\noindent\fbox{\parbox{\dimexpr\linewidth-2\fboxsep-2\fboxrule\relax}{\textbf{C4.}
[1.1 pt] From the graphs constructed in C3, determine $T_{eq}$. It is not
necessary to calculate uncertainty for this question.}}

If the variables are measured a finite number of times, Equation (4) can be
expressed as
\begin{equation}
\frac{T_{IC,j} - T_{IC,(j-1)}}{t_j - t_{j-1}} \propto \bar{T}_{OC} - \bar{T}_{IC} \tag{5}
\end{equation}
where $j$ represents the measurement number and $\overline{T}$ are the average
temperatures at each cylinder during the interval $t_{j-1}$ to $t_j$

\noindent\fbox{\parbox{\dimexpr\linewidth-2\fboxsep-2\fboxrule\relax}{\textbf{C5.}
[1.0 pt] Graph the expression on the left hand side of Equation 5 as a function
of the expression on the right hand side}}

\noindent\fbox{\parbox{\dimexpr\linewidth-2\fboxsep-2\fboxrule\relax}{\textbf{C6.}
[1.6 pt] Determine the value of $R_{Th}$ from the graph constructed in C5.}}

Fourier’s law for radial heat conduction through a slim cylindrical wall can be
written as:
\begin{equation}
\tfrac{dQ}{dt} = -\lambda A \tfrac{dT}{dr} \tag{6}
\end{equation}
where $A$ is the area of the wall, $\uplambda$ is the thermal conductivity of the
wall material, and r denotes the distance from the axis of the cylinder.

\noindent\fbox{\parbox{\dimexpr\linewidth-2\fboxsep-2\fboxrule\relax}{\textbf{C7.}
[1.6 pt] Using equations (4) and (6), determine the value $\uplambda$ of the
thermal conductivity of acrylic (the material separating IC and OC), indicating
the formula that you used.}}
